\section{Full-Rate Waveform Ingestion}
\label{sec:fullrate}

\subsection{Purpose}
The JSON path (Sections~5 and 6) stores one 20-sample cycle per channel every 5~s, a small fraction of the approximately 30~kHz the nodes acquire. To store every sample, nodes built with \texttt{MQTT\_PUBLISH\_DAQ2} publish each acquisition packet in binary on \texttt{daq2/<node\_id>/samples} (see the MQTT design document), and a dedicated service, \texttt{daq2-ingest}, writes them to InfluxDB.

\subsection{Why a Dedicated Service}
Telegraf cannot be used for this stream. Each DAQ2 packet carries 64 frames whose individual sample times must be computed from the header's first-frame time, frame period, and per-channel offsets, which no Telegraf parser supports. At the full rate, two nodes also produce approximately 120\,000 points per second, beyond what a JSON-based pipeline would sustain efficiently. \texttt{daq2-ingest} is a small Python service (\path{daq2-ingest/}) running as a fifth container in the Docker Compose stack. It subscribes to \texttt{daq2/+/samples}, decodes each packet, and writes gzip-compressed line protocol to the InfluxDB HTTP API in batches.

\subsection{Schema}
\begin{table}[H]
\centering
\caption{Full-Rate Schema (all points tagged \texttt{node\_id})}
\begin{tabular}{p{0.14\linewidth}p{0.14\linewidth}p{0.38\linewidth}p{0.22\linewidth}}
\hline
\textbf{Bucket} & \textbf{Measurement} & \textbf{Fields} & \textbf{One point per} \\
\hline
\texttt{waveform} & \texttt{voltage} & \texttt{a}, \texttt{b}, \texttt{c} (integer mV) & frame (\textasciitilde 33~\textmu s) \\
 & \texttt{current} & \texttt{a}, \texttt{b}, \texttt{c} (integer mA) & frame \\
 & \texttt{neutral} & \texttt{v}, \texttt{i} (neutral mode only) & frame \\
 & \texttt{daq2\_packet} & sequence, frame index, lock and time-valid flags, PTP and configuration epochs, timing uncertainty & packet (64 frames) \\
\texttt{waveform\_1s} & \texttt{voltage}, \texttt{current}, \texttt{neutral} & RMS, mean, minimum, and maximum of each field (float V/A), frame count, lock flag & UTC second \\
 & \texttt{power} & Real power, apparent power, and power factor per phase and in total, frame count, lock flag & UTC second \\
 & \texttt{daq2\_ingest} & service statistics & 10~s \\
\hline
\end{tabular}
\end{table}

Design decisions:
\begin{itemize}[noitemsep]
  \item \textbf{One point per signal group per frame.} The three phase voltages are sampled simultaneously in one ADC slot, and the three currents in the next, approximately 16.6~\textmu s later. Each group is written as one point at its own exact sampling time, rather than collapsing all six channels onto one timestamp, which would introduce a phase error of 0.3\textdegree{} at 50~Hz between voltage and current.
  \item \textbf{Integer storage.} Raw values are written as integers in mV and mA, which preserves the same 0.001~V/A precision as a three-decimal float. InfluxDB delta-encodes integers, and consecutive samples at 30~kHz differ only slightly, so integers compress approximately 4.8 times better than floats on noisy waveform data (Section~\ref{sec:sizing}).
  \item \textbf{Only trustworthy timestamps are written.} The firmware clears a packet's time-valid flag whenever the PTP servo adjusted the clock rate during it, which happens on every Sync (about once a second), as well as across clock steps and failed clock reads. A rate adjustment skews sample times by only nanoseconds over one packet, so such packets are kept if their first-frame time lies within 10~\textmu s of the time predicted from the previous packet and the clock epoch is unchanged; they are recorded with \texttt{time\_valid} = 0 in \texttt{daq2\_packet}. Packets spanning a clock step or showing a timestamp discontinuity are skipped and counted, as are packets more than a day from the server's clock, which a node publishes before its first PTP step while its clock still counts from boot. Before this rule, one 64-frame packet per second (0.2\% of the data) was discarded.
  \item \textbf{Timescale.} The nodes' TAI timestamps are converted to UTC with the offset the node received in the grandmaster's Announce messages, which each packet carries in its flags, as on the JSON path. Packets without one (before the node's first Announce) use \texttt{TAI\_UTC\_OFFSET\_S}, 37~s by default.
\end{itemize}

\subsection{Aggregation}
As it ingests each packet, the service accumulates per-channel sums for the current UTC second and emits the RMS, mean, minimum, and maximum once the second completes. It also computes power per phase following the IEEE~1459-2010 definitions: real power $P = \overline{v \cdot i}$, apparent power $S = V_\mathrm{rms} I_\mathrm{rms}$, and power factor $P/S$, with totals as the sums over the phases. Each ADC samples a phase's voltage and then its current one scan slot later (16.64~\textmu s, or 0.3\textdegree{} at 50~Hz), which uncorrected would understate real power by about 0.9\% at a power factor of 0.5. Each frame's current is therefore interpolated linearly to its voltage's sampling instant, between that frame's and the previous frame's current samples, using the slot offsets in the packet header; across a gap in the frame sequence, the first frame is excluded rather than interpolated. Unit tests confirm the computed power matches theory to within 0.05\% at power factors of 1 and 0.5. All of these are written to \texttt{waveform\_1s}, which is retained indefinitely, so that long-term trends remain available after the raw data expires. A complete second contains approximately 30\,048 frames; a lower \texttt{frames} count marks a partial second at start-up or after a gap. This replaces the Flux downsampling task of the original design (Section~\ref{sec:downsampling}).

\subsection{Process Structure and Flow Control}
Python executes one thread at a time per process, and a single-process design was measured to saturate at approximately 2.8 nodes, at which point the broker silently discarded the excess QoS~0 messages for every node. The service therefore runs as one receiver process, which only reads MQTT, and several worker processes (by default one per CPU core, at most eight). The receiver assigns each newly seen node to the worker with the fewest nodes, and all of that node's packets then go to the same worker, which holds its gap-detection, timestamp-continuity, and aggregation state. Each worker decodes, builds line protocol, and posts batches to InfluxDB over four parallel connections, since each post mostly waits on InfluxDB. If InfluxDB falls behind, up to approximately 10~s of raw data is buffered in memory before the oldest raw lines are dropped; aggregates are never dropped. If a worker's queue fills, the receiver drops that node's packets and counts them. Every loss is counted rather than silent: each process logs, and writes to \texttt{daq2\_ingest} (tagged by worker) every 10~s, the number of packets and frames received, gaps in the nodes' frame indices (data acquired but not delivered), packets skipped for invalid timestamps, queue and buffer drops, and write errors.

\subsection{Storage Sizing}
\label{sec:sizing}
Storage was measured by writing 60~s of synthetic three-phase data for one node, with Gaussian noise added (0.3~V and 0.03~A RMS) so that it compresses no better than a real signal:

\begin{table}[H]
\centering
\caption{Measured Storage per Value}
\begin{tabular}{lll}
\hline
\textbf{Format} & \textbf{Bytes per value} & \textbf{Two nodes, per day} \\
\hline
Float (V/A, 3 decimals) & 7.9 & \textasciitilde 246~GB \\
Integer (mV/mA) & 1.63 & \textasciitilde 51~GB \\
\hline
\end{tabular}
\end{table}

Integer storage is therefore approximately 4.8 times smaller, and is the default. At approximately 51~GB per day for two nodes, the default 7-day raw retention requires approximately 360~GB, and 30 days approximately 1.5~TB; requirements scale linearly with the number of nodes. The \texttt{waveform\_1s} bucket grows by approximately 1~MB per day. Raw retention should be set from the historian server's available disk, leaving headroom for InfluxDB compactions, and checked against actual usage once real data is flowing.

\subsection{Verification}
The service was verified on a development machine (16 cores, 16~GB RAM) in three ways:
\begin{itemize}[noitemsep]
  \item \textbf{Unit tests} check packet decoding against the firmware's \texttt{daq\_packet\_t} layout, sample time calculation, per-group timestamps, neutral-mode packets, rejection of malformed packets, integer conversion, gap and invalid-time counting, and per-second RMS (a 230~V sine yields 230.0~V).
  \item \textbf{A load test} published synthetic packets for two nodes at the firmware's real rate (30\,048 frames/s each) for 60~s. All 56\,342 packets were received, and InfluxDB held exactly 1\,802\,944 voltage and 1\,802\,944 current points per node, with no drops, gaps, or write errors. The service used approximately 55\% of one CPU core and InfluxDB approximately 77\%, with writes occupying approximately 42\% of the writer's time.
  \item The per-second aggregates reported 230.0~V RMS and 30\,048 frames for each complete second.
\end{itemize}
On hardware, a node built with \texttt{MQTT\_PUBLISH\_DAQ2} sustained approximately 470 packets/s (6.3~Mbit/s) with no drops at the node, and the service stored all 30\,048 frames per second.

Capacity was measured on a home test server (12 threads, 7.7~GB RAM) with one real node plus synthetic nodes:
\begin{table}[H]
\centering
\caption{Measured Capacity}
\begin{tabular}{lll}
\hline
\textbf{Nodes at full rate} & \textbf{Result} & \textbf{CPU used} \\
\hline
1 & No loss & ingest \textasciitilde 0.5 cores, InfluxDB \textasciitilde 0.5 \\
5 & No loss over 2 minutes; queues flat & ingest \textasciitilde 3.2 cores, InfluxDB \textasciitilde 7.8 \\
7 & \textasciitilde 9\% of packets dropped (CPU exhausted) & load average 16 on 12 threads \\
\hline
\end{tabular}
\end{table}
Capacity is therefore bounded by the server's CPU, at approximately 1.5 cores of InfluxDB and 0.6 cores of ingestion per full-rate node, and by disk, at approximately 25~GB per node per day.
